\begin{abstract}
Multiple-choice scores are commonly interpreted against chance. We argue that this comparison is often insufficient: before an arm is compared with another arm, its score should be tested against the construct's best constant, input-blind predictor. This floor is a necessary, not sufficient, validity condition. Across ten target constructs and one negative control, the floor is not chance and, for content scoring, is under-specified by the tie convention, length unit, and tokenizer; when option count varies, even ``chance'' is ambiguous. The floor must itself be calibrated: because it is a maximum over $k$ noisy label marginals, it is upward biased even under exactly uniform labels, and the calibrating null must be realisable by the construct --- a test we initially failed on our own largest benchmark. Under a null restricted to each item's legal letters, only two of the eight letter constructs (MMLU, BoolQ, both fixed-$k$) have floors a balanced null could not produce ($p<10^{-5}$); on the remaining six the floor lies inside the estimator's own sampling noise ($p=0.083$--$0.853$), so those constructs establish that chance is the wrong \emph{reference} without establishing that the floor differs from it in magnitude. On MMLU-Pro, all 21 evaluated cells are powered at the scale of the reference effect. Among 17 designated damaged cells in four model families, the null choice alone reverses the reading: damaged arms appear above a chance line while only two clear their best-constant floor, and the honest aggregate is 15/17 at or below the floor. Held to a single evidentiary standard --- a paired bootstrap interval excluding zero on \emph{both} sides, not only the floor side --- the flip is 3/12 above chance versus 1/12 above the floor, about a third of the size the asymmetric point-estimate comparison suggests. Across five smaller benchmarks, the designated set gives 9/85 cells above their floor against 45/85 above chance, the same wrong-null flip, while a mandatory power analysis explains why most per-benchmark significance tests are inconclusive. We then adapt a complementary, arm-conditional permutation null for a different question: whether one arm's predictions contain item-level information. Its statistic is the numerator of Cohen's $\kappa$ evaluated within option-count strata, so it is exactly zero for every pure constant emitter independent of collapse letter; our contribution is the stratification and its use as a pre-comparison gate, not the statistic. We use it to re-judge 27 existing cells. The re-analysis withdraws one above-floor capability label, dissolves both below-floor competence labels, and exposes a weak intact-family anchor. Finally, full-fp32 evaluation removes every bf16 exact tie and changes 18.03\% of letter decisions without improving accuracy, ruling out a numerical-tie explanation for the measurement failure. Together, the results motivate a simple reporting rule: calibrate each reported construct to its explicit input-blind null before interpreting arm differences.
\end{abstract}
